\documentclass[12pt,a4paper]{article}
\usepackage[utf8]{vntex, inputenc}
\usepackage{amsmath, amssymb, anyfontsize, array, color,  enumerate, enumitem, fancyhdr, float, graphicx, hyperref, lastpage, listings, mathtools, mips, natbib, pifont, tikz}
\usepackage[open,openlevel=0]{bookmark}
\usepackage[framemethod=tikz]{mdframed}

\usetikzlibrary{calc}
\newcommand\HRule{\rule{12cm}{1pt}}

\pagestyle{fancy}
\fancyhead{} % clear all header fields
\fancyhead[L]{
	\color{blue}
	\begin{tabular}{rl}
		\begin{picture}(25,15)(0,0)
		\put(0,-8){\includegraphics[width=8mm, height=8mm]{BK.png}}
		\end{picture}
		\begin{tabular}{l}
			\textbf{\bf \ttfamily Ho Chi Minh University of Technology}\\
			\textbf{\bf \ttfamily Faculty of Computer Science \& Engineering}
		\end{tabular}
	\end{tabular}
}
\fancyhead[R]{
	\begin{tabular}{l}
		\tiny \bf \\
		\tiny \bf
	\end{tabular}
}
\fancyfoot{} % clear all footer fields
\fancyfoot[L]{\scriptsize \ttfamily Computer Architecture Week 2}
\fancyfoot[R]{\scriptsize \ttfamily Page {\thepage}/\pageref{LastPage}}
\renewcommand{\headrulewidth}{0.3pt}
\renewcommand{\footrulewidth}{0.3pt}

\definecolor{dkgreen}{rgb}{0,0.6,0}
\definecolor{gray}{rgb}{0.5,0.5,0.5}
\definecolor{mauve}{rgb}{0.58,0,0.82}

\lstset{frame=tb,
	language=[MIPS]Assembler,
	aboveskip=3mm,
	belowskip=3mm,
	showstringspaces=false,
	columns=flexible,
	basicstyle={\small\ttfamily},
	numbers=none,
	numberstyle=\tiny\color{gray},
	keywordstyle=\color{blue},
	commentstyle=\color{dkgreen},
	stringstyle=\color{mauve},
	breaklines=true,
	breakatwhitespace=true,
	tabsize=3,
	numbers=left,
	stepnumber=1,
	numbersep=1pt,
	firstnumber=1,
	numberfirstline=true
}

\begin{document}
\begin{titlepage}
  \begin{tikzpicture}[remember picture, overlay]
    \draw[line width = 2pt] ($(current page.north west) + (1in,-1in)$) rectangle ($(current page.south east) + (-1in,1in)$);
  \end{tikzpicture}

  \begin{center}

    % Upper part of the page
    \textbf{\fontsize{12pt}{1pt}\selectfont HO CHI MINH CITY UNIVERSITY OF TECHNOLOGY}\\[0.5cm]
    {\fontsize{13pt}{1pt}\selectfont Faculty of Computer Science \& Engineering}\\[0.5cm]
    \begin{figure}[H]
      \centering
      \includegraphics[width=1.7in,height=1.7in]{BK.png}
    \end{figure}

    % Title
    \HRule\\[0.5cm]
    { \textbf{\fontsize{25pt}{1pt}\selectfont Computer Architecture Week 2}}\\[0.4cm]

    \HRule\\[0.8cm]
    \begin{minipage}{0.545\textwidth}
      \begin{flushleft}
        \textbf{Authors:}\\
        \begin{tabular}{l l}
          Lưu Nguyễn Hoàng Minh & 1952845 \\
          Vũ Anh Nhi            & 1952380 \\
          Nguyễn Chính Khôi     & 1952793 \\
          Nguyễn Hoàng          & 1952255 \\
        \end{tabular}
      \end{flushleft}
    \end{minipage}
    \begin{minipage}{0.4\textwidth}
      \begin{flushright}
        \textbf{Class:}\\
        CO2008\_CC02\\
        \textbf{Lecturer:}\\
        Kiều Đỗ Nguyên Bình\\

      \end{flushright}
    \end{minipage}

    \vfill

    % Bottom of the page
    \vspace{2cm}
    {\large} %{\large \today}
  \end{center}
\end{titlepage}


\pdfbookmark[0]{Question 1}{ques1}
\section*{Question 1:}

Write a simple MIPS program that can execute these steps: \\
1. Print a sentence to terminal to request an integer number from user; \\
2. Collect the number and increase it by 1;\\
3. Print the result to terminal.


\begin{mdframed}[hidealllines=true,backgroundcolor=magenta!10]
  \begin{lstlisting}
    # Declare main as a global function
  .text

  .globl main

  main:
    # Print string msg1
    li	$v0,4		# print_string syscall code = 4
    la	$a0, msg1	# load the address of msg1
    syscall

    # Get input A from user and save
    li	$v0,5		# read_int syscall code = 5
    syscall
    move	$t0,$v0		# syscall results returned in $v0

    # Main
    addi	$t0, $t0, 1	# A = A + 1

    # Print string msg3
    li	$v0, 4
    la	$a0, msg2
    syscall

    # Print sum
    li	$v0,1		# print_int syscall code = 1
    move	$a0, $t0	# int to print must be loaded into $a0
    syscall

    li	$v0,10		# exit
    syscall

    # Start .data segment (data!)
    .data
  msg1:	.asciiz	"Enter an integer: "
  msg2:	.asciiz	"New result is: "
  newline:   .asciiz	"\n"

          \end{lstlisting}
\end{mdframed}

\pdfbookmark[0]{Question 2}{ques2}
\section*{Question 2:}Write a small program that is able collect two integer numbers from users and print out the
sum of the two numbers

\begin{mdframed}[hidealllines=true,backgroundcolor=magenta!10]
  \begin{lstlisting}
    # Start .text segment (program code)
    .text

    .globl	main
  main:
    # Print string msg1
    li	$v0,4		# print_string syscall code = 4
    la	$a0, msg1	# load the address of msg
    syscall

    # Get input A from user and save
    li	$v0,5		# read_int syscall code = 5
    syscall
    move	$t0,$v0		# syscall results returned in $v0

    # Print string msg2
    li	$v0,4		# print_string syscall code = 4
    la	$a0, msg2	# load the address of msg2
    syscall

    # Get input B from user and save
    li	$v0,5		# read_int syscall code = 5
    syscall
    move	$t1,$v0		# syscall results returned in $v0

    # Math!
    add	$t0, $t0, $t1	# A = A + B

    # Print string msg3
    li	$v0, 4
    la	$a0, msg3
    syscall

    # Print sum
    li	$v0,1		# print_int syscall code = 1
    move	$a0, $t0	# int to print must be loaded into $a0
    syscall

    # Print \n
    li	$v0,4		# print_string syscall code = 4
    la	$a0, newline
    syscall

    li	$v0,10		# exit
    syscall

    # Start .data segment (data!)
    .data
  msg1:	.asciiz	"Enter A:   "
  msg2:	.asciiz	"Enter B:   "
  msg3:	.asciiz	"A + B = "
  newline:   .asciiz	"\n"

          \end{lstlisting}
\end{mdframed}

\pdfbookmark[0]{Question 3}{ques3}
\section*{Question 3:} Write a small program that allows users to input values for variables a, b, c, and d. The
program then calculates the following expressions and prints the results to terminal
\begin{equation}
  f = ( a + b ) - ( c - d - 2 )
\end{equation}
\begin{equation}
  g = ( a + b )*3 - ( c + d )*2
\end{equation}
\begin{mdframed}[hidealllines=true,backgroundcolor=magenta!10]
  \begin{lstlisting}
    .text

    .globl	main
  main:
    # Print string msg1
    li	$v0,4		# print_string syscall code = 4
    la	$a0, msg1	# load the address of msg
    syscall
    # Get input A from user and save
    li	$v0,5		# read_int syscall code = 5
    syscall
    move	$t0,$v0		# syscall results returned in $v0

    # Print string msg2
    li	$v0,4		# print_string syscall code = 4
    la	$a0, msg2	# load the address of msg2
    syscall
    # Get input B from user and save
    li	$v0,5		# read_int syscall code = 5
    syscall
    move	$t1,$v0		# syscall results returned in $v0

    # Print string msg3
    li	$v0,4		# print_string syscall code = 4
    la	$a0, msg3	# load the address of msg
    syscall
    # Get input C from user and save
    li	$v0,5		# read_int syscall code = 5
    syscall
    move	$t2,$v0		# syscall results returned in $v0

    # Print string msg4
    li	$v0,4		# print_string syscall code = 4
    la	$a0, msg4	# load the address of msg
    syscall
    # Get input D from user and save
    li	$v0,5		# read_int syscall code = 5
    syscall
    move	$t3,$v0		# syscall results returned in $v0




    # Calculate f = (a + b) - (c - d - 2)
    add	$t4, $t0, $t1	# t4 = a + b
    sub	$t5, $t2, $t3	# t5 = c - d
    addi	$t5, $t5, -2	# t5 = c - d - 2
    sub	$s0, $t4, $t5
    # Done

    # Calculate g = (a + b) * 3 - (c + d) * 2
    add $t6, $t4, $t4	# t6 = (a+b)*2
    add $t6, $t6, $t4	# t6 = (a+6)*3
    add $t7, $t2, $t3	# t7 = c+d
    add $t7, $t7, $t7	# t7 = (c+d)*2
    sub $s1, $t6, $t7



    # Print string msg5
    li	$v0, 4
    la	$a0, msg5
    syscall

    # Print sum
    li	$v0,1		# print_int syscall code = 1
    move	$a0, $s0	# int to print must be loaded into $a0
    syscall

    # Print \n
    li	$v0,4		# print_string syscall code = 4
    la	$a0, newline
    syscall


    # Print string msg6
    li	$v0, 4
    la	$a0, msg6
    syscall

    # Print sum
    li	$v0,1		# print_int syscall code = 1
    move	$a0, $s1	# int to print must be loaded into $a0
    syscall

    # Print \n
    li	$v0,4		# print_string syscall code = 4
    la	$a0, newline
    syscall

    li	$v0,10		# exit
    syscall

    # Start .data segment (data!)
    .data
  msg1:	.asciiz	"Enter A:   "
  msg2:	.asciiz	"Enter B:   "
  msg3:	.asciiz	"Enter C:   "
  msg4:	.asciiz	"Enter D:   "

  msg5:	.asciiz	"f = (a + b) - (c - d - 2) = "
  msg6:	.asciiz	"g = (a + b) * 3 - (c + d) * 2 = "

  newline:   .asciiz	"\n"

        \end{lstlisting}
\end{mdframed}


\pdfbookmark[0]{Question 4}{ques4}
\section*{Question 4:}

Write a small program that allows users to input 5 different integer numbers. The program then prints those numbers in reverse. For example, users input 1, 2, 3, 4, 5; the program should print 5, 4, 3, 2, 1.

\begin{mdframed}[hidealllines=true,backgroundcolor=magenta!10]
  \begin{lstlisting}

    .text

    .globl	main
  main:
    # Print string msg1
    li	$v0,4
    la	$a0, msg1
    syscall
    # Get input A from user and save
    li	$v0,5
    syscall
    move	$t0,$v0

    # Print string msg2
    li	$v0,4
    la	$a0, msg2
    syscall
    # Get input B from user and save
    li	$v0,5
    syscall
    move	$t1,$v0

    # Print string msg3
    li	$v0,4
    syscall
    # Get input C from user and save
    li	$v0,5
    syscall
    move	$t2,$v0


    # Print string msg4
    li	$v0,4
    la	$a0, msg4
    syscall
    # Get input D from user and save
    li	$v0,5
    syscall
    move	$t3,$v0


    # Print string msg5
    li	$v0,4
    la	$a0, msg5
    syscall

    # Get input E from user and save
    li	$v0,5
    syscall
    move	$t4,$v0

    # Print string msg6
    li	$v0, 4
    la	$a0, msg6
    syscall

    # Print sum
    li	$v0,1
    move	$a0, $t4
    syscall
    # Print ' '
    li	$v0,4
    la	$a0, space
    syscall
    # Print sum
    li	$v0,1
    move	$a0, $t3
    syscall
    # Print ' '
    li	$v0,4
    la	$a0, space
    syscall
    # Print sum
    li	$v0,1
    move	$a0, $t2
    syscall
    # Print ' '
    li	$v0,4
    la	$a0, space
    syscall
    # Print sum
    li	$v0,1
    move	$a0, $t1
    syscall
    # Print ' '
    li	$v0,4
    la	$a0, space
    syscall
    # Print sum
    li	$v0,1
    move	$a0, $t0
    syscall
    # Print ' '
    li	$v0,4
    la	$a0, space
    syscall
    # Print \n
    li	$v0,4
    la	$a0, newline
    syscall

    li	$v0,10		# exit
    syscall

    # Start .data segment (data!)
    .data
  msg1:	.asciiz	"Enter A:   "
  msg2:	.asciiz	"Enter B:   "
  msg3:	.asciiz	"Enter C:   "
  msg4:	.asciiz	"Enter D:   "
  msg5:	.asciiz	"Enter E:   "
  msg6:	.asciiz	"Print in reverse order:      "
  newline:   .asciiz	"\n"
  space:  .asciiz     "  "


          \end{lstlisting}
\end{mdframed}

\pdfbookmark[0]{Question 5}{ques5}
\section*{Question 5:}

Write a MIPS program that calculates and prints the result of the following equation to the terminal with the number of instructions as small as possible.

\(66000+30+30+30+30+30+30+30+30+30+30 - 6000+25\)


\begin{mdframed}[hidealllines=true,backgroundcolor=magenta!10]
  \begin{lstlisting}

    # Declare main as a global function
    .globl main
    .text

  # The label 'main' represents the starting point
  main:
    li $t0, 30
    add $t0, $t0, $t0
    add $t1, $t0, $t0
    add $t1, $t1, $t1
    add $t1, $t1, $t0

    addi $t1, $t1, 66000
    addi $t1, $t1, -6000
    addi $a0, $t1, 25

    li $v0, 1 # Sets $v0 to "10" to select exit syscall
    syscall # Exit

          \end{lstlisting}
\end{mdframed}

\pdfbookmark[0]{Question 6}{ques6}
\section*{Question 6:}

Write a MIPS program that receives values for a, b, c, d, and x variables from users and prints
result of following equation:

\(f = a \times x^3 + b \times x^2 + c \times x + d\)


\begin{mdframed}[hidealllines=true,backgroundcolor=magenta!10]
  \begin{lstlisting}
  .text

  .globl main
    j main

  MULTS:
    #a * b = res. a = t0, b = t1, res = s0

    #int res = 0;
    addi $sp, $sp, -4
    sw $s0, 0($sp)
    addi $s0, $zero, 0

    #int i = 0; while (i != b) {res += a; i++;}
    addi $sp, $sp, -4	#int i = 0
    sw $s1, 0($sp)
    addi $s1, $zero, 0	#int i = 0 - END

  WHILE:	beq $s1, $t1, EXIT
    add $s0, $s0, $t0
    addi $s1, $s1, 1
    j WHILE

  EXIT:	move $v0, $s0 		#return
    lw $s1, 0($sp)		#restore s1
    addi $sp, $sp, 4
    lw $s0, 0($sp)		#restore s0
    addi $sp, $sp, 4

    jr $ra

  main:
    # x = s0, a = s1, b = s2, c = s3, d = s4, result = s5

    # Get x
    li	$v0,4
    la	$a0, msg1
    syscall
    li	$v0,5
    syscall
    move	$s0,$v0

    # Get a
    li	$v0,4
    la	$a0, msg2
    syscall
    li	$v0,5
    syscall
    move	$s1,$v0

    # Get b
    li	$v0,4
    la	$a0, msg3
    syscall
    li	$v0,5
    syscall
    move	$s2,$v0

    # Get c
    li	$v0,4
    la	$a0, msg4
    syscall
    li	$v0,5
    syscall
    move	$s3,$v0

    # Get d
    li	$v0,4
    la	$a0, msg5
    syscall
    li	$v0,5
    syscall
    move	$s4,$v0

    # Main Program
    addi $s5, $zero, 0 	#res = 0;
    addi $t2, $s0, 0	#temp = x;

    add $s5, $s5, $s4	#res += d;

    addi $t0, $s0, 0	#res += x * c;
    addi $t1, $s3, 0
    jal MULTS
    add $s5, $s5, $v0

    addi $t0, $s0, 0	#x = x * temp;
    addi $t1, $t2, 0
    jal MULTS
    addi $s0, $v0, 0

    addi $t0, $s0, 0	#res += x * b;
    addi $t1, $s2, 0
    jal MULTS
    add $s5, $s5, $v0

    addi $t0, $s0, 0	#x = x * temp;
    addi $t1, $t2, 0
    jal MULTS
    addi $s0, $v0, 0

    addi $t0, $s0, 0	#res += x * a;
    addi $t1, $s1, 0
    jal MULTS
    add $s5, $s5, $v0

    # Print
    li	$v0,4
    la	$a0, msg6
    syscall
    li	$v0,1
    move	$a0, $s5
    syscall
    li	$v0,10
    syscall

    .data
  msg1:	.asciiz	"Enter x: "
  msg2:	.asciiz	"Enter a: "
  msg3:	.asciiz	"Enter b: "
  msg4:	.asciiz	"Enter c: "
  msg5:	.asciiz	"Enter d: "
  msg6:	.asciiz	"Result is: "
  newline:   .asciiz	"\n"

          \end{lstlisting}
\end{mdframed}
\end{document}
